\section{Energy Context and Prediction}
\label{sec:energy}

\noindent\textbf{Evaluation scope.} This section defines how NRE and prediction enter the MCS in simulation. Fig.~\ref{fig:eval-energy} reports internal NRE traces (instrumentation check), not Energest-calibrated joules or duty-cycle measurements.

Sections~\ref{sec:arch} and~\ref{sec:sec} established how QoS and trust indicators enter the MCS. This section clarifies how energy enters the same score, distinguishes \emph{software-simulated} energy dynamics in Cooja from calibrated field harvesting, and documents the bounded memory footprint of the predictor~\cite{osterlind2006cooja,oikonomou2022contik}.

\subsection{Normalized residual energy (NRE)}
The reference module exposes the normalized residual energy indicator (NRE) as an eight-bit value on the $[0,100]$ scale through \texttt{aer\_energy\_get\_nre\_x100()}, derived from an internal fixed-point register \texttt{residual\_x100} updated on each call to \texttt{aer\_rpl\_energy\_periodic()}. NRE therefore tracks a \emph{logical} state variable suitable for coupling to the MCS via \texttt{aer\_rpl\_plus\_update\_context()}, not a coulomb-counted battery model tied to a specific sensor board.

\subsection{Prediction as a short-horizon moving average}
The ``predicted'' energy term returned by \texttt{aer\_energy\_get\_pred\_x100()} is the mean of the last \texttt{N\_HISTORY} samples (eight in the current source) of \texttt{residual\_x100}, again scaled to $[0,100]$. This is a lightweight surrogate for heavier predictors (e.g., recurrent neural networks) discussed in prior AER-RPL prototypes~\cite{belacel2025aerrpl}: it offers smoothing and trend visibility with \emph{O}(1) update cost and static RAM.

\subsection{Cooja-only evolution of the residual state}
\texttt{aer\_rpl\_energy\_periodic()} adjusts \texttt{residual\_x100} by subtracting a pseudo-random drain term and adding a pseudo-random harvest term within documented bounds, clamping the register to remain inside $[500,10000]$ in fixed-point hundredths. \textbf{This process models abstract charge variability for simulation experiments; it is not calibrated to photovoltaic traces, fuel-gauge ICs, or outdoor irradiance.} Accordingly, Section~\ref{sec:eval} exposes only internal NRE traces (Fig.~\ref{fig:eval-energy}) as an instrumentation check---never as validated millijoule or duty-cycle measurements from a physical testbed.

\subsection{Interaction with context fusion \texorpdfstring{$(\alpha,\beta,\gamma)$}{alpha/beta/gamma}}
Low NRE reduces the context coefficient $\gamma$ inside \texttt{aer\_rpl\_plus\_update\_context()}, increasing $\beta$ and thereby emphasizing the energy partial score within the MCS (Section~\ref{sec:arch}). High NRE produces a mild opposite adjustment to $\gamma$. This coupling keeps energy stress \emph{visible} at parent-selection time rather than relegating it to offline post-processing.

\subsection{Relation to analytical energy models in the literature}
Analytical and simulation studies of energy-aware RPL provide useful baselines for understanding how rank dynamics interact with duty cycling~\cite{limjoco2018analytical}. AER-MQoS complements such work empirically through Cooja sweeps with open scripts; future work may tighten the link between the abstract drain/harvest process and published energy models, but only after the additional parameters are versioned alongside the datasets.

\subsection{Summary}
The energy plane contributes two bounded indicators---NRE and a moving-average prediction---to the MCS through a transparent Cooja-oriented process, consistent with the simulation-only evaluation scope stated in Section~\ref{sec:ql}.
